%************************************************
\chapter{Tool Implementation}\label{ch:rlrom}
%************************************************
We implemented our methods into RLRom \footnote{\url{https://github.com/decyphir/rlrom}}, part of the impremenation was reported in the paper \cite{gomezruiz2025overlay}, which integrates a rich STL parser with efficient online monitoring methods 
(taken from \cite{donze2013,deshmukh_robust_2015})  with popular reinforcement learning frameworks. Our tool is similar to stl-gym\footnote{\url{https://github.com/nphamilton/stl-gym}} 
based on \cite{nickovic_rtamt_2020}, which allows to replace the reward function with the robustness of an STL formula. RLRom goes beyond by allowing richer formulas to be written and monitored, 
as well as to train models for specifications described by reward machines with STL predicates.  

RLRom is built over three main existing components, namely:
\begin{itemize}
\item \textbf{Gymnasium} \cite{Towers2024Gymnasium}: A Python library offering a standardized API for developing and comparing reinforcement learning algorithms by providing a variety 
of pre-built and customizable environments. It serves as a widely used platform for reinforcement learning research and development.
\item \textbf{Stable-Baselines3} \cite{stable-baselines3}: A Python library that provides well-documented implementations of model-free reinforcement learning algorithms 
(such as PPO, DQN, SAC, etc) in PyTorch. It offers a consistent interface for training and evaluating reinforcement learning agents.
\item \textbf{STLRom\footnote{\url{https://github.com/decyphir/stlrom}}}: A C++ library with Python bindings for robust online monitoring of Signal Temporal Logic (STL). It enables users to 
define temporal properties of signals and compute their robustness value online.
\end{itemize}

The interface with Gymnasium is based on the wrapper mechanism, which takes an existing reinforcement learning environment with a reset(), step() and reward() methods and observation and action 
spaces, and returns a new environment that changes one or all of these components. To interface with STLRom, we implemented an stl\_wrapper that does the following:
\begin{enumerate}
    \item It adds an stl\_driver object from STLRom which can read data and monitors STL formulas.
    \item Within the step() function, a new observation of the wrapped environment is added to the stl\_driver and the satisfaction of the formulas is updated; the resulting robustness values 
    are then used to augment the observation and/or influence the reward.
\end{enumerate}

\section{Workflows and features}
\label{sec:workflows}

When working with a new reinforcement learning problem, assuming we are given a reinforcement learning environment implementing a Gymnasium-compatible API, RLRom is intended to support the following steps, with workflows of increasing 
complexities.
\begin{enumerate}
\item \textbf{(WF1) Testing}: RLRom provides a simple way to run multiple episodes varying initial seed. Models can be imported from files or from popular repos, such as Hugging Face 
\footnote{\url{https://huggingface.com}} Eventually, we plan on implementing more advanced simulation-based testing methods, in the spirit of \cite{donze_breach_2010}. 
\item \textbf{(WF2) Training}: new models can be trained with RLRom interfacing to existing stable-baselines3 algorithms - existing hyper-parameters can be quickly fetched and 
tried to obtain some baseline model for the base problem;
\item \textbf{(WF3) Monitoring}: Once a baseline model has been obtained, STL specifications can be written and monitored. This requires deciding which feature of the observation space 
should be monitored as a signal, and what properties should be checked. A typical first step is to formalize the initial goal of the environment, and check that the baseline model satisfy 
this formalization. E.g., for the classic Cart-Pole environment, the goal is to keep the pole up, which can be translated into the simple predicate \texttt{pole\_up := abs(theta[t]) < epsi} 
for some small value of \texttt{epsi}. The goal specification can then be:
  \begin{center}
\begin{verbatim}
phi_goal_cartpole := ev_[0,400] alw_[0,100] mu
\end{verbatim}
  \end{center}
  that is, before 400 steps, the pole should stay up for at least 100 steps.
\item \textbf{(WF4) Training with RM-STL specifications}: Reward machines with STL can be incorporated into the model training to solve a problem with more complex objectives. 
\end{enumerate}

(WF1) and (WF2) are standard workflows in the reinforcement learning community, implemented for example in the RL-Zoo framework \footnote{\url{https://rlzoo.readthedocs.io}}. (WF3) and (WF4) 
are more unique to RLRom, which enabled us to obtain interesting results to complex reinforcement learning problems, described in the previous chapters. 

\section{Overall Architecture}

The workflows of Section~\ref{sec:workflows}, from testing (WF1) to
STL-conditioned reward machines (WF4), are all enabled by a common
architecture with the following design principles:

\begin{itemize}
\item \textbf{Configuration-driven.} Everything that varies between
  experiments (environment, algorithm, the STL specification etc)
  lives in a single YAML file. The Python code is generic and does not
  need to be touched to change an experiment.
\item \textbf{Wrapper-based integration.} The monitoring and reward
  machine logic is implemented as a stack of Gymnasium \emph{wrappers}
  inserted between the bare environment and the agent. This keeps the
  RL algorithm untouched and makes the machinery reusable across
  environments and algorithms.
\item \textbf{Separation of concerns and modularity} Sections in the
  configuration each correspond to a different mechanism or task
  (interface, test, train, STL, reward machines) and can be
  interchanged and re-used between experiments.
\end{itemize}

A run is described by a single YAML file with a small set of \emph{top-level} keys,
summarized in Table~\ref{tab:config}. Sub-values can be factored out: values ending
in \texttt{.yml} are recursively replaced by the content of the referenced
file, and a \texttt{.stl} field loads the
text of an STL specification file.

\begin{table}[htbp]
  \centering
  \caption{Main top-level keys of a \textsc{RLROM} configuration file.}
  \label{tab:config}
  \small
  \begin{tabular}{@{}llp{9.2cm}@{}}
    \toprule
    \textbf{Key} &  \textbf{Role} \\
    \midrule
    \texttt{env\_name}  & environment (e.g., Gymnasium id) \\
    \texttt{model\_path}, \texttt{model\_name} & where to read/write model checkpoints \\
    \texttt{model\_use\_specs}  & whether the model was trained with the wrapped
                                  observation/reward\\
                                  &(allows to monitor pre-trained models) \\
    \texttt{cfg\_specs}  & the STL and reward machines specification and how it is used \\
    &(signals,
      reward/eval/label formulas, reward machines declaration) \\
    \texttt{cfg\_test}  & evaluation: seeds, number of episodes/steps, rendering \\
    \texttt{cfg\_train}  & training: algorithm, timesteps, eval cadence, repetition \\
    \bottomrule
  \end{tabular}
\end{table}

Listing~\ref{lst:min} (\texttt{cfg0\_min.yml}) is a minimal 
configuration: an environment and a test block. 

\begin{cfglisting}[caption={Minimal configuration
(\texttt{examples/cartpole/cfg0\_min.yml}).},label={lst:min}]
env_name: CartPole-v1
cfg_test:
  init_seeds: 1
  num_ep: 5
  num_steps: 100
  render_mode: 'human'
\end{cfglisting}

This will test the environment without any agent (actions are then chosen randomly).

\section{Testing and training standard agents (WF1 and WF2)}
\label{sec:testing}

The first two workflows are the ``standard'' ones that RLRom shares with the
broader reinforcement-learning ecosystem (they are, for instance, at the core of
the RL-Zoo framework): \textbf{(WF1)} testing an existing agent and
\textbf{(WF2)} training a new one. Both are illustrated here on the CartPole
environment used as the running example of this chapter. Thanks to the
configuration-driven design of Section~\ref{sec:design:config}, switching between
them --- and between environments, algorithms and models --- requires no change to
the Python code, only an edit to the YAML.

\subsection{Testing an agent imported from HuggingFace (WF1)}

The simplest way to obtain a ready-made agent is to import a pretrained model from
a public repository such as HuggingFace\footnote{\url{https://huggingface.co}}.
Listing~\ref{lst:hugimport} shows \texttt{cfg0\_hug.yml}, which loads the PPO
checkpoint \texttt{sb3/ppo-CartPole-v1} for CartPole: setting
\texttt{model\_path} to \texttt{huggingface} tells RLRom to fetch the model from
HuggingFace under the name \texttt{sb3/ppo-CartPole-v1}. The \texttt{cfg\_test}
block describes five episodes of at most $100$ steps rendered to screen.

\begin{cfglisting}[caption={Testing a HuggingFace agent
(\texttt{examples/cartpole/cfg0\_hug.yml}).},label={lst:hugimport}]
env_name: CartPole-v1
model_path: huggingface
model_name: sb3/ppo-CartPole-v1

cfg_test:
  init_seeds: 0
  num_ep: 5
  num_steps: 100
  render_mode: 'human'
\end{cfglisting}

\noindent
The run is launched with
\begin{center}\ttfamily rlr test examples/cartpole/cfg0\_hug.yml\end{center}
which loads the checkpoint and plays the requested episodes, logging the per-
episode reward and a summary of the runs. Choosing \texttt{model\_path: random}
(resp.\ \texttt{keyboard}) instead substitutes a random (resp.\ human) policy,
which is a convenient sanity check of the environment's dynamics before any real
evaluation.

\subsection{Training a baseline agent (WF2)}

To train a new model, RLRom interfaces with the algorithms implemented in
Stable-Baselines3 and keeps all training options in a configuration file.
Listing~\ref{lst:trppo} shows \texttt{cfg0tr\_ppo.yml}, which trains a PPO agent
for the bare CartPole task (no STL specification yet): \texttt{cfg\_test}
describes the evaluation, and \texttt{cfg\_train} selects the algorithm by name
(\texttt{ppo}), pulls its hyper-parameters from the referenced
\texttt{cartpole\_ppo.yml} (Listing~\ref{lst:ppohp}), and sets the number of
parallel environments, the total number of timesteps, the evaluation frequency and
how many times the whole training is repeated from scratch
(\texttt{repeat\_training}, used later to estimate variance across seeds).

\begin{center}
\begin{minipage}[t]{0.47\textwidth}
\centering
\begin{cfglisting}[caption={Baseline training
(\texttt{examples/cartpole/cfg0tr\_ppo.yml}).},label={lst:trppo}]
env_name: CartPole-v1
model_path: ./models
model_name: vanilla_ppo

cfg_test:
  init_seeds: 0
  num_ep: 10
  num_steps: 200
  render_mode: 'human'

cfg_train:
  n_envs: 8
  total_timesteps: 100_000
  eval_freq: 10_000
  repeat_training: 5
  algo:
    ppo: cartpole_ppo.yml
\end{cfglisting}
\end{minipage}\hfill
\begin{minipage}[t]{0.47\textwidth}
\centering
\begin{cfglisting}[caption={PPO hyper-parameters
(\texttt{examples/cartpole/cartpole\_ppo.yml}).},label={lst:ppohp}]
device: cpu
verbose: 0

policy_kwargs:
  net_arch: [32, 32]

n_steps: 1024
learning_rate: 3e-4
\end{cfglisting}
\end{minipage}
\end{center}

\noindent
Training is launched with
\begin{center}\ttfamily rlr train examples/cartpole/cfg0tr\_ppo.yml\end{center}
which instantiates $8$ parallel environments, runs PPO for $100\,000$ timesteps
and repeats the experiment $5$ times from scratch. Figure~\ref{fig:vanilla} shows
the resulting learning curves: Figure~\ref{fig:vanilla:raw} plots the mean episode
reward of each of the $5$ runs, and Figure~\ref{fig:vanilla:env} overlays, in
addition, the min--max band across runs (the ``envelope'' produced by
\texttt{rlr plot --envelope}). The band gives an immediate sense of the variance
between seeds, which is why the rest of this chapter reports training with the
envelope view.

\begin{figure}[htbp]
  \centering
  \subfloat[All \texttt{vanilla\_ppo} trainings.]{
    \includegraphics[width=0.92\textwidth]
      {../examples/cartpole/vanilla_ppo_trainings.png}
    \label{fig:vanilla:raw}
  }\\[1em]
  \subfloat[Envelope (mean and min--max band) across trainings.]{
    \includegraphics[width=0.92\textwidth]
      {../examples/cartpole/vanilla_ppo_trainings_enveloppe.png}
    \label{fig:vanilla:env}
  }
  \caption{Learning curves of the \texttt{vanilla\_ppo} baseline (WF2), from the
  checkpoint result files produced by \texttt{rlr train}.}
  \label{fig:vanilla}
\end{figure}

\paragraph{Other algorithms.}
Although this chapter uses PPO throughout, the integration is generic: any
algorithm implemented in Stable-Baselines3 (DQN, SAC, A2C, ...) can be selected by
simply changing \texttt{cfg\_train.algo} to the corresponding key (e.g.,
\texttt{dqn: cartpole\_dqn.yml}) and pointing at a matching hyper-parameter file.
The wrapper-based design keeps the agent \emph{agnostic} to the STL machinery: the
same algorithms -- and largely the same hyper-parameters as in
Listing~\ref{lst:ppohp} -- can be kept for the specification-aware workflows that
follow, which isolates the effect of adding the STL signal from the effect of
switching the optimization algorithm.

%%% Local Variables:
%%% mode: LaTeX
%%% TeX-master: "../rlrom_chapter"
%%% End:
